% =====================================================================
% CITAÇÕES E REFERÊNCIAS
%
% No texto você usa \cite{chave} ou \citeonline{chave}; as obras ficam
% cadastradas em bibliografia.bib. Nunca escreva a lista de referências
% à mão: ela é montada a partir das citações.
%
% As opções abaixo são do pacote abntex2cite e definem a aparência:
%   alf ............ sistema autor-data, "(SILVA, 2020)". A alternativa
%                    seria o sistema numérico, "[1]".
%   versalete ...... imprime os sobrenomes em VERSALETE (letras maiúsculas
%                    de tamanho reduzido), como em Silva.
%   abnt-emphasize=bf .......... destaca o título em negrito
%   abnt-etal-list=3 ........... a partir de 3 autores, usa "et al."
%   abnt-etal-text=it .......... "et al." em itálico
%   abnt-and-type=& ............ liga os autores com "&"
%   abnt-last-names=abnt ....... formata os sobrenomes no padrão ABNT
%   abnt-repeated-author-omit=no  repete o nome do autor a cada entrada,
%                                 em vez de substituí-lo por um traço
%
% Confirme com seu curso qual sistema ele exige antes de mudar.
% =====================================================================
\usepackage[
  alf,
  versalete,
  abnt-emphasize=bf,
  abnt-etal-list=3,
  abnt-etal-text=it,
  abnt-and-type=&,
  abnt-last-names=abnt,
  abnt-repeated-author-omit=no
]{abntex2cite}

% Texto do link que leva da referência de volta à página em que foi citada.
\renewcommand{\backref}{}
\renewcommand*{\backrefalt}[4]{%
  \ifcase #1
    Nenhuma citação no texto.%
  \or
    Citado na página #2.%
  \else
    Citado nas páginas #2.%
  \fi}
